\begin{examplebox}
	\textbf{\textcolor{CExample}{例 1（基础）}}
	\textbf{\textcolor{CExample}{题目：}}
	已知点 $P$ 与 $P'$ 关于直线 $l$ 对称，$PP'$ 交 $l$ 于点 $M$，且 $PP' = 8$。求 $PM$ 的长度。
	\textbf{\textcolor{CExample}{解题步骤：}}
	\begin{description}[style=nextline, leftmargin=2.8em, labelsep=0.5em, itemsep=0.45em]
		\item[\textbf{第一步（应用性质）}]
			由对称性质，$l$ 是 $PP'$ 的垂直平分线，所以 $M$ 是 $PP'$ 的中点。
			\par\textbf{理由：} 轴对称基本性质。
		\item[\textbf{第二步（计算长度）}]
			$PM = \dfrac{1}{2}\,PP' = \dfrac{1}{2}\times 8 = 4$。
			\par\textbf{理由：} 中点等分线段。
		\item[\textbf{第三步（写出答案）}]
			所以 $PM = 4$。
			\par\textbf{理由：} 由上述计算。
	\end{description}
	\textbf{\textcolor{CExample}{关键结论：}}
	\textbf{$PM = 4$。}

	\medskip
	\textbf{\textcolor{CExample}{例 2（稍变形）}}
	\textbf{\textcolor{CExample}{题目：}}
	在平面直角坐标系中，已知 $A(2,0)$，$B(0,2)$。判断直线 $l: x+y-2=0$ 是否为线段 $AB$ 的垂直平分线？$A$ 与 $B$ 是否关于 $l$ 对称？
	\textbf{\textcolor{CExample}{解题步骤：}}
	\begin{description}[style=nextline, leftmargin=2.8em, labelsep=0.5em, itemsep=0.45em]
		\item[\textbf{第一步（求中点）}]
			$AB$ 中点 $M$ 的坐标为 $\left(\dfrac{2+0}{2},\dfrac{0+2}{2}\right) = (1,1)$。
			\par\textbf{理由：} 中点坐标公式。
		\item[\textbf{第二步（验证点在线上）}]
			把 $M$ 代入 $l$，$1+1-2=0$，满足，所以 $l$ 经过 $AB$ 的中点。
			\par\textbf{理由：} 中垂线必过中点。
		\item[\textbf{第三步（验证垂直）}]
			$k_{AB} = \dfrac{2-0}{0-2} = -1$，$k_l = -1$。由于 $k_{AB}\cdot k_l = 1 \neq -1$，故 $l$ 与 $AB$ 不垂直。因此 $l$ 不是 $AB$ 的中垂线，$A$ 与 $B$ 不关于 $l$ 对称。
			\par\textbf{理由：} 中垂线必须垂直于线段。
	\end{description}
	\textbf{\textcolor{CExample}{关键结论：}}
	\textbf{ $l$ 不是 $AB$ 的中垂线，$A$ 与 $B$ 不关于 $l$ 对称。}

	\medskip
	\textbf{\textcolor{CExample}{例 3（综合应用）}}
	\textbf{\textcolor{CExample}{题目：}}
	已知点 $A(0,2)$，求 $A$ 关于直线 $l: 2x-y-1=0$ 的对称点 $A'$ 的坐标。
	\textbf{\textcolor{CExample}{解题步骤：}}
	\begin{description}[style=nextline, leftmargin=2.8em, labelsep=0.5em, itemsep=0.45em]
		\item[\textbf{第一步（设点列中点方程）}]
			设 $A'(x,y)$，则 $AA'$ 中点 $\left(\dfrac{x}{2},\dfrac{y+2}{2}\right)$ 在 $l$ 上：$\;2\cdot\dfrac{x}{2} - \dfrac{y+2}{2} -1 = 0$，化简得 $2x - y - 4 = 0$。
			\par\textbf{理由：} 对称点连线中点在对称轴上。
		\item[\textbf{第二步（列垂直方程）}]
			由 $AA'\perp l$，斜率之积为 $-1$：$\dfrac{y-2}{x}\cdot 2 = -1$，整理得 $x+2y-4=0$。
			\par\textbf{理由：} 垂直条件。
		\item[\textbf{第三步（解方程组得坐标）}]
			联立 $
			\begin{cases}
				2x - y -4 =0 \\ x+2y-4 =0
		\end{cases}
			$，解得 $x=\frac{12}{5},\ y=\frac{4}{5}$。所以 $A'(\frac{12}{5},\frac{4}{5})$。
			\par\textbf{理由：} 解线性方程组。
	\end{description}
	\textbf{\textcolor{CExample}{关键结论：}}
	\textbf{ $A'(\frac{12}{5},\frac{4}{5})$。}
\end{examplebox}
